\documentclass[a4paper,12pt]{article}
\usepackage[utf8]{inputenc}
\usepackage[T1]{fontenc}
\usepackage[ngerman]{babel}
\usepackage{geometry}
\usepackage{hyperref}
\geometry{margin=2.5cm}

\begin{document}

\section*{Physics-based Machine Learning for Mantle Convection Simulations}
\textbf{Autoren:} Siddhant Agarwal, Ali Can Bekar, Christian H\"uttig,
David S.\ Greenberg, Nicola Tosi \\
\textbf{Erscheinungsjahr:} 2025 (arXiv:2505.16041) \\
\textbf{Institution:} DLR Berlin \& Helmholtz-Zentrum Hereon, Geesthacht

\subsection*{Problemstellung und Motivation}

Mantelkonvektionssimulationen sind ein zentrales Werkzeug zur Untersuchung
der thermischen Evolution von Gesteinsplaneten wie Erde, Mars und Venus.
Die zugrundeliegenden partiellen Differentialgleichungen -- Massen-, Impuls-
und Energieerhaltung -- beschreiben ein tr\"ages, hochviskoses Fluid im sog.
\emph{Stokes-Regime} ($Re \ll 1$), bei dem Tr\"agheitskr\"afte vernachl\"assigbar
sind. Der numerisch aufwendigste Schritt jeder Zeitintegration ist dabei die
L\"osung des Stokes-Problems (Massen- und Impulserhaltung), das einen direkten
Gleichungsl\"oser erfordert und den Gro{\ss}teil der Rechenzeit verbraucht.
Erschwerend kommen schlecht bekannte Eingangsparameter sowie stark nichtlineare
Viskosit\"atsabh\"angigkeiten von Druck und Temperatur hinzu, die ausgedehnte
Parameterstudien erforderlich machen. Die Autoren stellen daher die Frage, ob
ein physik-basiertes ML-Modell den Stokes-L\"oser ersetzen und damit die
Simulation erheblich beschleunigen kann.

\subsection*{Methodik}

\paragraph{Hybrides Simulations-Framework.}
Der Kernansatz ist eine hybride Kopplung: Ein trainiertes konvolutionales
neuronales Netz (CNN) ersetzt den Stokes-L\"oser und sagt aus dem Temperaturfeld
$T^t$ die Str\"omungsgeschwindigkeiten $(u^t, v^t)$ vorher. Diese werden
anschlie{\ss}end an den numerischen Finite-Volumen-L\"oser GAIA \"ubergeben,
der daraus das Temperaturfeld $T^{t+1}$ per Advektions-Diffusions-Schritt
bestimmt. Auf diese Weise wird die Zeitintegration \emph{ohne} zeitliches
Lernen des Netzes durchgef\"uhrt, was die physikalische Kausali\"at der Stokes-
Gleichungen ausnutzt: Die Geschwindigkeit zu einem gegebenen Zeitschritt h\"angt
nicht von der Geschwindigkeit des vorherigen Zeitschritts ab, sondern
ausschlie{\ss}lich vom aktuellen Temperatur- und Viskosit\"atsfeld.

\paragraph{CNN-Architektur.}
Die Architektur folgt Tompson et al.\ (2017) und verwendet f\"unf Aufl\"osungsebenen
mit Average-Pooling und bikubischem Upsampling. Im Gegensatz zum klassischen
U-Net werden die Merkmale der Encoder- und Decoderst\"ufen durch
Konkatenation statt durch Skip-Connections summiert, was dem Netz mehr
Flexibilit\"at bei der Kombination von Information unterschiedlicher Skalen
gibt. Als Aktivierungsfunktion wird GELU mit Group Normalisation eingesetzt.
Das Netz verarbeitet neben dem Temperaturfeld auch das Viskosit\"atsfeld sowie
Gitterkoordinaten und Simulationsparameter als Eingaben.

\paragraph{Massenerhaltung durch Vektorpotential (Stromfunktion).}
Um Divergenzfreiheit der vorhergesagten Geschwindigkeitsfelder
mathematisch zu gew\"ahrleisten, sagt das Netz nicht $(u, v)$ direkt
vorher, sondern ein skalares Vektorpotential $a$ (die
zweidimensionale Stromfunktion). Die divergenzfreien Geschwindigkeitskomponenten
ergeben sich dann durch:
\[
  u = \frac{\partial a}{\partial y}, \qquad v = -\frac{\partial a}{\partial x},
\]
wobei die Ableitungen \"uber feste Finite-Differenzen-Faltungskerne berechnet
werden. Dieser \emph{Hard Constraint} garantiert Massenerhaltung im Innern
des Gebiets bis auf Maschinengenauigkeit (double precision).

\paragraph{Gelernte Randauff\"ullungen.}
Standardm\"a{\ss}ige Padding-Strategien (Zero- oder Replicate-Padding) f\"uhren
an den Gebietsstr\"andern zu erheblichen Fehlern, die sich im autoregressiven
Rollout aufw\"urfen und die L\"osung destabilisieren. Die Autoren setzen
stattdessen \emph{Boundary Learned Convolutions} ein (Innamorati et al.\ 2020):
Ein separater Satz von Faltungsfiltern wird f\"ur jede der 8 Randregionen
(4 Kanten, 4 Ecken) sowie das Innere trainiert. Dies reduziert den MAE an
den R\"andern um bis zu Faktor~8 gegen\"uber Zero-Padding.

\paragraph{Geschwindigkeitsskalierung und Verlustfunktion.}
Die Trainingsdaten umfassen Str\"omungsgeschwindigkeiten, die \"uber mehrere
Gr\"o{\ss}enordnungen variieren -- je nach Simulationsparametern ($Q$, $\gamma$,
$\beta$). Standardnormierungen wie Min-Max-Skalierung versagen hier.
Stattdessen wird ein parametrisch angepasster Skalierungsfaktor
\[
  s = 5\,e^{0.18\,Q}\,\gamma^{0.43}\,\beta^{-0.46}
\]
genutzt, der die Geschwindigkeitsfelder verschiedener Simulationen in
einen homogeneren Bereich bringt. Die Verlustfunktion kombiniert MAE
auf Geschwindigkeiten und deren Ableitungen, einen Randstrafterm sowie
eine dynamische Verlustgewichtung, die kleinen Geschwindigkeitsamplituden
mehr Aufmerksamkeit schenkt.

\subsection*{Ergebnisse}

Mit nur 94 Trainingssimulationen erreicht das Modell eine bis zu
\textbf{89-fache Beschleunigung} gegen\"uber dem direkten CPU-L\"oser und
\"ubertrifft damit auch alternative Baselines (Momentum-Skips, iterativer
GPU-L\"oser) sowohl in Genauigkeit als auch in Geschwindigkeit.
Stabile autoreg\-ressive Rollouts \"uber Zehntausende von Zeitschritten
sind in nahezu allen Testf\"allen m\"oglich. Im Gegensatz dazu divergiert
ein reines U-Net, das direkt in der Zeit lernt (\emph{Learning in Time}),
bereits nach 16 Schritten -- trotz besserer Einzel-Schritt-Vorhersagegenauigkeit.
Dies unterstreicht, dass physikalische Strukturen (Stokes-Kausalit\"at)
fundamentaler f\"ur die Rollout-Stabilit\"at sind als eine niedrige
Validierungs-Loss. Schw\"achen zeigen sich bei Extrapolation im Parameterraum
sowie bei stark au{\ss}erhalb der Trainingsverteilung liegenden
Anfangsbedingungen (\emph{Out-of-Distribution}-F\"alle).

\subsection*{Bezug zur Masterarbeit}

Diese Arbeit ist in mehrfacher Hinsicht unmittelbar relevant f\"ur die
Masterarbeit zur physik-basierten ML-Vorhersage von Str\"omungsfeldern
um sedimentierende Kristalle.

\paragraph{Gemeinsames physikalisches Regime.}
Beide Arbeiten operieren im Stokes-Str\"omungsregime ($Re \ll 1$), in dem
Tr\"agheitskr\"afte vernachl\"assigt werden k\"onnen und das Stokes-Problem
den Rechenengpass darstellt. Die hier entwickelten methodischen L\"osungen
sind daher direkt \"ubertragbar.

\paragraph{Stromfunktions-Ansatz zur Gew\"ahrleistung von Divergenzfreiheit.}
Der Vektorpotential-Ansatz in Gl.\ (5) entspricht exakt der Stromfunktions-
Formulierung ($v = \nabla \times \psi$), die in der Masterarbeit als Alternative
zur direkten Geschwindigkeitsvorhersage untersucht wird. Diese Arbeit
best\"atigt empirisch, dass der Hard Constraint deutlich robustere
Rollout-Eigenschaften liefert als ein Soft Constraint in der Verlustfunktion --
ein wichtiger Befund f\"ur die Designentscheidung der Masterarbeit.

\paragraph{Residuelles Lernen und physikalische Vorkenntnisse.}
Der hybride Ansatz -- Netz lernt Korrektur zum physikalisch motivierten
Hintergrundfeld -- ist konzeptuell verwandt mit dem in der Masterarbeit
implementierten Residuallernen ($v_\text{total} = v_\text{Stokes} + \Delta v_\text{gelernt}$).
Beide Ans\"atze nutzen analytische L\"osungen als starke induktive Vorannahme
und lernen nur die verbleibende Abweichung, was den Lernproblem erheblich
vereinfacht.

\paragraph{Herausforderung der Geschwindigkeitsskalierung.}
Die parametrisch motivierte Skalierungsstrategie in Gl.\ (4) adressiert
dasselbe Problem, das auch in der Masterarbeit auftritt: Geschwindigkeitsfelder
verschiedener Konfigurationen (hier: Simulationsparameter; dort: Kristallanzahl
und -positionen) k\"onnen stark in ihrer Amplitude variieren. Die hier
vorgeschlagene Regressionsanpassung an Systemparameter bietet eine
\"ubertragbare Strategie f\"ur die Normierung in der Masterarbeit.

\paragraph{Vergleich UNet vs.\ CNN-Architekturen.}
Die systematische Gegen\"uberstellung von reinem U-Net (Learning in Time)
und dem Stokes-Surrogat-CNN liefert wichtige Argumente f\"ur die
Architekturwahl in der Masterarbeit. Die \"Uberlegenheit des physikalisch
strukturierten Ansatzes -- trotz schlechterer Einzelschritt-Genauigkeit --
unterst\"utzt die These, dass physikalische Nebenbedingungen wichtiger sind
als Modellkapazit\"at allein.

\paragraph{Gelernte Randbehandlung.}
Die starke Sensitivit\"at des autoregressiven Rollouts gegen\"uber
Randfehlern und der deutliche Vorteil gelernter Randauff\"ullungen
(Faktor~8 genauer) sind auch f\"ur die UNet-basierte Kristallsedimentations-
vorhersage relevant, da Randeffekte um einzelne Kristalle \"ahnliche
Instabilit\"aten hervorrufen k\"onnen.

\paragraph{Grenzen der Generalisierung.}
Das klare Scheitern des Modells bei Out-of-Distribution-Bedingungen (neue
Energiegleichung, kalte Startwerte) spiegelt die zentrale Forschungsfrage der
Masterarbeit wider: Kann ein auf $N$ Kristalle trainiertes Modell auf
$N + m$ Kristalle generalisieren? Diese Arbeit zeigt, dass auch physik-
informierte ML-Modelle empfindlich auf Verteilungsverschiebungen reagieren
und betont die Notwendigkeit eines repr\"asentativen Trainingsdatensatzes --
ein Argument f\"ur die in der Masterarbeit implementierte Datenerweiterung
durch Spiegelungen und Translationen.

\end{document}